\documentclass[aspectratio=169]{beamer}
\usepackage[utf8]{inputenc}
\usepackage[T1]{fontenc}
\usepackage[es-ES]{babel}
\usepackage{graphicx}
\usepackage{tikz}
\usepackage{listings}
\usepackage{xcolor}

% Tema
\usetheme{Madrid}
\usecolortheme{default}

% Colores personalizados
\definecolor{unsazul}{RGB}{30, 60, 120}
\definecolor{unsverde}{RGB}{0, 100, 0}

\setbeamercolor{structure}{fg=unsazul}
\setbeamercolor{palette primary}{bg=unsazul,fg=white}
\setbeamercolor{palette secondary}{bg=unsazul!80,fg=white}
\setbeamercolor{palette tertiary}{bg=unsazul!60,fg=white}

% Configuración de listings
\lstset{
    basicstyle=\ttfamily\small,
    keywordstyle=\color{unsazul}\bfseries,
    commentstyle=\color{gray}\itshape,
    stringstyle=\color{unsverde},
    showstringspaces=false,
    breaklines=true,
    frame=single,
    backgroundcolor=\color{gray!10}
}

% Información del documento
\title{Métodos de Análisis de Datos 1}
\subtitle{Clase 5: Transformación de Datos}
\author{Universidad Nacional del Sur}
\date{}

\begin{document}

% Portada
\begin{frame}
\titlepage
\end{frame}

% Contenidos
\begin{frame}{Contenidos de hoy}
\tableofcontents
\end{frame}

\section{Conversión de tipos}

\begin{frame}{¿Por qué importan los tipos de datos?}
\begin{alertblock}{Problema común}
Los datos se cargan con tipos incorrectos
\end{alertblock}

\begin{exampleblock}{Ejemplos}
\begin{itemize}
    \item Números como texto: \texttt{"123"} en vez de \texttt{123}
    \item Fechas como texto: \texttt{"2024-01-15"} en vez de \texttt{datetime}
    \item Categóricas como texto: \texttt{"rojo"} en vez de \texttt{category}
    \item Booleanos como números: \texttt{1/0} en vez de \texttt{True/False}
\end{itemize}
\end{exampleblock}

\pause

\begin{block}{Consecuencias}
\begin{itemize}
    \item Operaciones incorrectas (\texttt{"1" + "2" = "12"})
    \item Más memoria usada
    \item Visualizaciones erróneas
    \item Modelos de ML fallan
\end{itemize}
\end{block}
\end{frame}

\begin{frame}{Tipos de datos principales}
\begin{columns}[t]
\column{0.5\textwidth}
\textbf{Numéricos:}
\begin{itemize}
    \item \texttt{int}: Enteros
    \item \texttt{float}: Decimales
\end{itemize}

\textbf{Texto:}
\begin{itemize}
    \item \texttt{string}: Texto general
    \item \texttt{category}: Categóricas (menos memoria)
\end{itemize}

\column{0.5\textwidth}
\textbf{Booleanos:}
\begin{itemize}
    \item \texttt{bool}: True/False
\end{itemize}

\textbf{Fechas:}
\begin{itemize}
    \item \texttt{datetime}: Fecha y hora
    \item \texttt{date}: Solo fecha
\end{itemize}

\textbf{Categóricos ordenados:}
\begin{itemize}
    \item \texttt{ordinal}: Con orden natural
\end{itemize}
\end{columns}
\end{frame}

\begin{frame}[fragile]{Conversión de tipos en Python}
\begin{lstlisting}[language=Python]
import pandas as pd

# A numérico (errores → NaN)
df['edad'] = pd.to_numeric(df['edad'], errors='coerce')

# A string
df['id'] = df['id'].astype(str)

# A fecha
df['fecha'] = pd.to_datetime(df['fecha'], format='%Y-%m-%d')

# A categórico
df['color'] = df['color'].astype('category')

# Categórico ordenado
df['educacion'] = pd.Categorical(
    df['educacion'],
    categories=['primaria', 'secundaria', 'universitaria'],
    ordered=True
)
\end{lstlisting}
\end{frame}

\begin{frame}[fragile]{Conversión de tipos en R}
\begin{lstlisting}[language=R]
library(lubridate)

# A numérico
df$edad <- as.numeric(df$edad)

# A string
df$id <- as.character(df$id)

# A fecha
df$fecha <- ymd(df$fecha)  # año-mes-día
# Otras: dmy(), mdy(), ymd_hms()

# A factor
df$color <- as.factor(df$color)

# Factor ordenado
df$educacion <- factor(
  df$educacion,
  levels = c('primaria', 'secundaria', 'universitaria'),
  ordered = TRUE
)
\end{lstlisting}
\end{frame}

\section{Variables derivadas}

\begin{frame}{Feature Engineering básico}
\begin{block}{¿Qué son variables derivadas?}
Nuevas variables creadas a partir de las existentes
\end{block}

\begin{block}{¿Por qué crearlas?}
\begin{itemize}
    \item Extraer información útil
    \item Capturar relaciones no lineales
    \item Mejorar performance de modelos
    \item Facilitar interpretación
\end{itemize}
\end{block}

\pause

\begin{exampleblock}{Ejemplos}
\begin{itemize}
    \item IMC = peso / altura²
    \item Edad = hoy - fecha\_nacimiento
    \item Nombre\_completo = nombre + apellido
    \item Es\_adulto = edad >= 18
    \item Trimestre = ceil(mes / 3)
\end{itemize}
\end{exampleblock}
\end{frame}

\begin{frame}[fragile]{Variables derivadas: Python}
\begin{lstlisting}[language=Python]
# Operaciones aritméticas
df['imc'] = df['peso'] / (df['altura'] ** 2)

# Operaciones con fechas
df['edad'] = (pd.Timestamp.now() - df['fecha_nac']).dt.days / 365.25

# Extraer componentes de fecha
df['año'] = df['fecha'].dt.year
df['mes'] = df['fecha'].dt.month
df['dia_semana'] = df['fecha'].dt.dayofweek  # 0=lunes

# Operaciones con texto
df['nombre_completo'] = df['nombre'] + ' ' + df['apellido']
df['longitud_nombre'] = df['nombre'].str.len()

# Condiciones
df['adulto'] = df['edad'] >= 18
df['categoria_edad'] = pd.cut(
    df['edad'],
    bins=[0, 18, 65, 100],
    labels=['joven', 'adulto', 'mayor']
)
\end{lstlisting}
\end{frame}

\begin{frame}[fragile]{Variables derivadas: R}
\begin{lstlisting}[language=R]
library(tidyverse)

df <- df %>%
  mutate(
    # Aritmética
    imc = peso / (altura^2),
    
    # Fechas
    edad = as.numeric(difftime(Sys.Date(), fecha_nac, units = "days")) / 365.25,
    año = year(fecha),
    mes = month(fecha),
    dia_semana = wday(fecha),  # 1=domingo
    
    # Texto
    nombre_completo = paste(nombre, apellido),
    longitud_nombre = nchar(nombre),
    
    # Condiciones
    adulto = edad >= 18,
    categoria_edad = cut(edad, 
                        breaks = c(0, 18, 65, 100),
                        labels = c('joven', 'adulto', 'mayor'))
  )
\end{lstlisting}
\end{frame}

\section{Codificación de categóricas}

\begin{frame}{El problema de las variables categóricas}
\begin{alertblock}{Algoritmos de ML requieren números}
La mayoría de algoritmos no pueden trabajar directamente con texto
\end{alertblock}

\begin{exampleblock}{Ejemplo}
\begin{center}
\begin{tabular}{|c|}
\hline
\textbf{Color} \\
\hline
rojo \\
verde \\
azul \\
\hline
\end{tabular}
\quad $\rightarrow$ \quad ¿Cómo convertir a números?
\end{center}
\end{exampleblock}

\pause

\begin{block}{Dos estrategias principales}
\begin{enumerate}
    \item \textbf{Label Encoding:} Asignar un número a cada categoría
    \item \textbf{One-Hot Encoding:} Crear columna binaria por categoría
\end{enumerate}
\end{block}
\end{frame}

\begin{frame}{Label Encoding}
\begin{block}{Concepto}
Asigna un número único a cada categoría

\begin{center}
\begin{tabular}{|c|c|}
\hline
\textbf{Color} & \textbf{Codificado} \\
\hline
rojo & 0 \\
verde & 1 \\
azul & 2 \\
\hline
\end{tabular}
\end{center}
\end{block}

\pause

\begin{alertblock}{⚠ Cuidado}
\textbf{Solo usar para variables ordinales} (con orden natural)

Para variables nominales, el modelo asumirá que 2 > 1 > 0, lo cual no tiene sentido para colores
\end{alertblock}

\begin{block}{¿Cuándo usar?}
Educación: primaria (0) < secundaria (1) < universitaria (2) ✓
\end{block}
\end{frame}

\begin{frame}[fragile]{Label Encoding: Código}
\begin{columns}[t]
\column{0.5\textwidth}
\textbf{Python:}
\begin{lstlisting}[language=Python]
from sklearn.preprocessing import LabelEncoder

# Label encoding
le = LabelEncoder()
df['educacion_encoded'] = le.fit_transform(
    df['educacion']
)

# Ver mapeo
mapeo = dict(zip(
    le.classes_, 
    le.transform(le.classes_)
))
print(mapeo)
# {'primaria': 0, 
#  'secundaria': 1, 
#  'universitaria': 2}
\end{lstlisting}

\column{0.5\textwidth}
\textbf{R:}
\begin{lstlisting}[language=R]
# Automático con factor
df$educacion_encoded <- as.numeric(
  factor(df$educacion)
)

# Con orden específico
df$educacion_encoded <- as.numeric(
  factor(
    df$educacion,
    levels = c('primaria', 
               'secundaria', 
               'universitaria')
  )
)
\end{lstlisting}
\end{columns}
\end{frame}

\begin{frame}{One-Hot Encoding}
\begin{block}{Concepto}
Crea una columna binaria (0/1) para cada categoría
\end{block}

\begin{center}
\begin{tabular}{|c|c|c|c|}
\hline
\textbf{Original} & \textbf{color\_rojo} & \textbf{color\_verde} & \textbf{color\_azul} \\
\hline
rojo & 1 & 0 & 0 \\
verde & 0 & 1 & 0 \\
azul & 0 & 0 & 1 \\
rojo & 1 & 0 & 0 \\
\hline
\end{tabular}
\end{center}

\pause

\begin{block}{Ventajas}
\begin{itemize}
    \item No asume orden entre categorías
    \item Funciona para variables nominales
    \item Fácil de interpretar
\end{itemize}
\end{block}

\begin{alertblock}{Desventaja}
Aumenta mucho la dimensionalidad si hay muchas categorías
\end{alertblock}
\end{frame}

\begin{frame}[fragile]{One-Hot Encoding: Python}
\begin{lstlisting}[language=Python]
import pandas as pd

# Método simple
df_encoded = pd.get_dummies(
    df, 
    columns=['color'], 
    prefix='color'
)

# Con sklearn (más control)
from sklearn.preprocessing import OneHotEncoder

encoder = OneHotEncoder(
    sparse=False,      # Retornar array denso
    drop='first'       # Eliminar primera categoría 
)                      # (evita multicolinealidad)

encoded = encoder.fit_transform(df[['color']])

# A DataFrame
df_encoded = pd.DataFrame(
    encoded,
    columns=encoder.get_feature_names_out(['color'])
)
\end{lstlisting}
\end{frame}

\begin{frame}[fragile]{One-Hot Encoding: R}
\begin{lstlisting}[language=R]
library(caret)

# Con caret
dummies <- dummyVars(" ~ color", data = df)
df_encoded <- predict(dummies, newdata = df)

# Con tidyverse
df_encoded <- df %>%
  mutate(value = 1) %>%
  pivot_wider(
    names_from = color, 
    values_from = value, 
    values_fill = 0,
    names_prefix = "color_"
  )
\end{lstlisting}
\end{frame}

\begin{frame}{¿Label o One-Hot?}
\begin{columns}[t]
\column{0.5\textwidth}
\textbf{Label Encoding:}
\begin{itemize}
    \item ✓ Variables ordinales
    \item ✓ Árboles de decisión
    \item ✓ Menos columnas
    \item ✗ Variables nominales
    \item ✗ Regresión lineal
\end{itemize}

\column{0.5\textwidth}
\textbf{One-Hot Encoding:}
\begin{itemize}
    \item ✓ Variables nominales
    \item ✓ Regresión lineal/logística
    \item ✓ Neural networks
    \item ✗ Muchas categorías
    \item ✗ Aumenta dimensionalidad
\end{itemize}
\end{columns}

\vspace{1em}

\begin{exampleblock}{Regla práctica}
\begin{itemize}
    \item Ordinal → Label Encoding
    \item Nominal con pocas categorías (<10) → One-Hot
    \item Nominal con muchas categorías → Target Encoding, Frequency Encoding, etc.
\end{itemize}
\end{exampleblock}
\end{frame}

\section{Resumen}

\begin{frame}{Pipeline de transformación}
\begin{enumerate}
    \item \textbf{Verificar tipos de datos}
    \begin{itemize}
        \item \texttt{df.dtypes} (Python) o \texttt{str(df)} (R)
    \end{itemize}
    
    \item \textbf{Convertir tipos incorrectos}
    \begin{itemize}
        \item Fechas como datetime
        \item Números como int/float
        \item Categóricas como category/factor
    \end{itemize}
    
    \item \textbf{Crear variables derivadas}
    \begin{itemize}
        \item Extraer info de fechas
        \item Combinar columnas
        \item Aplicar fórmulas
    \end{itemize}
    
    \item \textbf{Codificar categóricas}
    \begin{itemize}
        \item Label encoding para ordinales
        \item One-hot encoding para nominales
    \end{itemize}
\end{enumerate}
\end{frame}

\begin{frame}{Resumen de la clase}
\begin{itemize}
    \item Los \textbf{tipos de datos} deben ser correctos
    \item Conversión en Python: \texttt{astype()}, \texttt{pd.to\_numeric()}, \texttt{pd.to\_datetime()}
    \item Conversión en R: \texttt{as.numeric()}, \texttt{ymd()}, \texttt{factor()}
    \item Las \textbf{variables derivadas} extraen información útil
    \item \textbf{Label encoding:} solo para ordinales
    \item \textbf{One-hot encoding:} para nominales
    \item La elección depende del tipo de variable y del algoritmo
\end{itemize}

\vspace{1em}

\begin{block}{Próxima clase}
\textbf{Normalización, estandarización e integración de datos}
\end{block}
\end{frame}

\begin{frame}{Tarea}
\begin{enumerate}
    \item Usar el dataset limpio de la clase anterior
    \item Transformaciones a realizar:
    \begin{itemize}
        \item Verificar y corregir tipos de datos
        \item Crear al menos 3 variables derivadas útiles
        \item Si hay fechas, extraer componentes (año, mes, día semana)
        \item Identificar variables categóricas
        \item Codificar variables categóricas apropiadamente
    \end{itemize}
    \item Documentar en el notebook:
    \begin{itemize}
        \item Qué transformaciones hiciste
        \item Por qué elegiste label o one-hot
        \item Comparar dimensiones antes vs después
    \end{itemize}
    \item Subir a GitHub
\end{enumerate}
\end{frame}

\begin{frame}[standout]
\Huge ¿Preguntas?
\end{frame}

\end{document}
